\section{Introduction}
\label{sec:intro}

Postoperative atrial fibrillation (POAF) complicates roughly one quarter of isolated coronary artery bypass grafting (CABG) operations and is associated with perioperative stroke, prolonged intensive care, readmission, higher expenditure, and reduced long-term survival.\cite{aranki1996predictors,villareal2004,dobrev2019postoperative,caldonazo2023meta,gaudino2023ehj} Targeted preventive measures---including posterior left pericardiotomy,\cite{gaudino2021palacs} prophylactic amiodarone or colchicine, and intensified telemetry\cite{chatterjee2026sts}---reduce POAF risk, yet universal prophylaxis exposes lower-risk patients to avoidable toxicity. Accurate risk stratification \emph{before skin incision} is therefore essential for both clinical trial enrichment and individualized perioperative care.

Yet the cardiac-surgery machine-learning literature displays a stark paradox. Published POAF models advertise areas under the receiver-operating-characteristic curve (ROC-AUCs) from 0.58 to 0.98, whereas two independent PROBAST critical appraisals concluded that over 80\% carry high risk of bias and shrink from a median apparent AUC of 0.71 to a median external AUC of 0.61.\cite{fields2025poaf,lopeztuesta2026poaf} Apparent AUCs above 0.75 almost invariably reflect post-incision or intensive-care predictors,\cite{mathew2004multicenter,magee2007atrial,ma2025noaf,cao2025mimic4} mixed valve-and-CABG cohorts,\cite{mariscalco2014poaf,burgos2021comaf} single-center internal validation,\cite{cho2026poaf} non-registry waveform or multimodal streams,\cite{han2025aiecg,khalpey2026digital} or continuous telemetry rather than registry-coded treated AF.\cite{filardo2017telemetry,perezgrovas2023jaha}

Compounding this comparability gap, national surgical registries expose two deeper methodological pitfalls. First, in the Society of Thoracic Surgeons Adult Cardiac Surgery Database (STS ACSD),\cite{sts2014v281,shahian2018sts,obrien2018sts2} conditional skip logic leaves uncomplicated recoveries blank, early table columns conceal postoperative durations, and default final-year (calendar-2017) holdouts share 99.5\% of hospitals with training. Second, binary $\{0,1\}$ supervision forces models to fit noisy outcomes shaped heavily by unmeasured intraoperative perturbations,\cite{echahidi2008mechanisms,dobrev2019postoperative} trapping gradient-boosted trees and frontier tabular foundation models (TabPFN-3.5,\cite{hollmann2025tabpfn,jaeger2026tabpfn35} KumoTabular-Large v2,\cite{qu2026kumotabular} and TabFM v1.0.0\cite{kong2026tabfm}) inside a narrow band.

We audited all 368,132 isolated-CABG procedures across 1,091 U.S.\ surgical centers in STS ACSD v2.81 to answer three questions: (1)~what is the true pre-incision discrimination, calibration, clinical utility, and equity of bedside risk proxies, boosted trees, and tabular foundation models on strictly unseen hospitals? (2)~Can evolutionary continuous-target learning (\textbf{EvoTarget-FT}), re-engineered parallel genetic programming (\textbf{JAIN Ultra GP v2}), and their orthogonal \textbf{Dual-Paradigm Fusion} break the 50,000-patient binary ceiling without leakage? (3)~How does performance shift from \emph{cold-start unseen hospitals} (0\% overlap) to a controlled \emph{in-network seen-hospital} regime (all 1,091 centers in training) with and without train-only empirical-Bayes institutional calibration?
